\section{Border models and smooth homogeneous forms}
\label{sec:models}

We first identify the two border models precisely. All closures in this
section are over $\C$. In a finite-dimensional polynomial space, coefficientwise
convergence means ordinary convergence of the coefficients of every monomial.
The image of a polynomial map is constructible
\cite[Theorem~29.23.3, Tag~054K]{StacksProject}, and a complex constructible
set has the same Euclidean and Zariski closures. For the latter fact, decompose
the set into finitely many locally closed subsets: each contains a Zariski
open dense subset of its closure, which is also Euclidean dense. These facts
apply both to the polynomial map taking an affine matrix to its determinant
and to the action map defining a linear-group orbit. Thus coefficientwise
sequences can be used for either closure below.

\begin{lemma}[Affine substitution]
\label{lem:affine-substitution}
Let $f$ be a complex polynomial, let $A$ be an affine map into its variable
space, and let $a\ne0$ be a complex scalar. Then
\[
 \bdc(a f\circ A)\le \bdc(f).
\]
\end{lemma}
\begin{proof}
Fix a size $n$ and determinants $Q_j=\det M_j$ converging coefficientwise to
$f$, with all entries of $M_j$ affine-linear. The entries of $M_j\circ A$
are still affine-linear. Substitution by the fixed map $A$ is a continuous
linear map on spaces of polynomials of degree at most $n$, so
$Q_j\circ A\to f\circ A$. Multiplying one row of each substituted matrix
by $a$ gives the assertion. This argument imposes no bound on the matrix
entries in the approximating sequence.
\end{proof}

For $n\ge m$, let $W_n=\Mat_n(\C)$ and regard $\det_n$ as a degree-$n$
form on $W_n$. Choose $m^2$ independent linear coordinates $x_{ij}$ on $W_n$.
If $n>m$, choose one more independent coordinate $z$ and set
\[
 P_{n,m}=z^{n-m}\per_m(x).
\]
If $n=m$, set $P_{m,m}=\per_m(x)$; no padding coordinate is needed.
Let $\mathcal D_n$ be the projective Zariski closure of the orbit of
$[\det_n]$ under invertible linear substitutions on $W_n$.

\begin{proposition}[Equivalence of the two border models]
\label{prop:model-equivalence}
For every $n\ge m\ge1$,
\[
 \bdc(\per_m)\le n
 \quad\Longleftrightarrow\quad
 [P_{n,m}]\in\mathcal D_n.
\]
\end{proposition}

\begin{proof}
Every linear map $L:W_n\to W_n$ is a limit of invertible ones. Consequently
$\det_n\circ L$ belongs to the affine closure of the determinant orbit.
This orbit is invariant under nonzero scalar multiplication: multiplying
one row of a matrix variable by $a\ne0$ multiplies its determinant by $a$
and is an invertible linear substitution. Its affine closure is therefore
the cone over $\mathcal D_n$. In particular a nonzero form belongs to this
affine closure exactly when its projective class belongs to $\mathcal D_n$.

Suppose first that $\bdc(\per_m)\le n$. Identity blocks let us take every
approximating matrix to have size exactly $n$. Write
\[
 Q_j(x)=\det\left(A_{j,0}+\sum_{a,b=1}^m x_{ab}A_{j,ab}\right)
 \longrightarrow\per_m(x).
\]
For $n>m$, homogenize to degree $n$ using the independent coordinate $z$:
\[
 Q_j^{\mathrm{hom}}(x,z)=\det\left(zA_{j,0}+\sum_{a,b=1}^m x_{ab}A_{j,ab}\right)
 \longrightarrow z^{n-m}\per_m(x).
\]
Homogenization here is a fixed linear map of coefficient spaces. Each
displayed determinant is a linear substitution of $\det_n$ on $W_n$,
so its limit belongs to the same closed cone. The coordinates fit because
$m^2+1\le n^2$ when $n>m$.

For $n=m$, instead take the degree-$m$ homogeneous component of each $Q_j$.
It is precisely
\[
 \det\left(\sum_{a,b=1}^m x_{ab}A_{j,ab}\right),
\]
and these components converge to $\per_m$. They again are linear
substitutions of $\det_m$. This proves the forward implication without
introducing an extra variable in the case $n=m$.

Conversely, the cone observation and equality of the two closures give
determinant-orbit representatives converging coefficientwise to $P_{n,m}$.
When $n>m$, set $z=1$ and all unused coordinates to zero; when $n=m$, just
use the $m^2$ coordinates $x_{ab}$. Each resulting polynomial is an affine
determinant of size $n$, and the limit is $\per_m$. This proves the converse.
Finally, a determinant of size less than $m$ has degree less than $m$, as
does every polynomial in its coefficientwise closure. Thus
$\bdc(\per_m)\ge m$, and the stated range $n\ge m$ covers every size
that can satisfy the affine condition.
\end{proof}

The quantitative obstruction used in the proof depends on the degree and
the number of variables of a smooth homogeneous polynomial. Here smoothness
means that the projective zero set has no point at which all first partial
derivatives vanish.

\begin{theorem}[Lower bound for smooth homogeneous forms]
\label{thm:smooth-border}
Let $G\in\C[x_1,\ldots,x_d]$ be a nonzero homogeneous polynomial of degree
$r\ge2$, where $d\ge2$. If its zero set in $\PP^{d-1}$ is smooth, then
\[
 \bdc(G)\ge \frac{(r-1)(d-1)}{4e}.
\]
\end{theorem}

The complete proof is in Section~\ref{sec:polar}. For a constant vector
$w\in\C^d$, write
\[
 \partial_wG=\sum_{j=1}^d w_j\frac{\partial G}{\partial x_j}
\]
for the first directional derivative, or \emph{first polar}, of $G$.
If $G$ is a coefficientwise limit of affine determinants of size $n$, the proof
chooses $d-2$ first polars whose common intersection with the hypersurface
consists of $r(r-1)^{d-2}$ simple points. It then shows that these selected
points persist for one sufficiently close determinant and satisfy
\begin{equation}
 r(r-1)^{d-2}
 \le 2^{d-2}\binom{2n-1}{d-1}
 \le \left(\frac{4en}{d-1}\right)^{d-1}.
 \label{eq:smooth-count-synopsis}
\end{equation}
The binomial coefficient is understood to be zero when $d-1>2n-1$;
in that case the positive left side already rules out the approximation.
The first inequality comes from lifting the selected points to the right
and left kernel lines of the determinant matrix and counting the resulting
bilinear equations. Taking the $(d-1)$st root gives the theorem. The proof
establishes persistence and simplicity of those selected lifts using only
coefficientwise closeness of the determinant polynomials, so it allows
unbounded coefficients in the approximating matrices.

This comparison specifies our construction's target. A nonzero homogeneous
form of degree proportional to $k$ in quadratically many variables, with
a smooth projective zero set, forces cubic determinant size. We now produce
such a form as a fixed affine restriction of a permanent whose size is
linear in $k$.
